\section{映射列的收敛}
\label{sec:mmap-map-sequences}

用多项式、截断级数或迭代格式逼近一个映射时，得到的是一列映射$f_1,f_2,\ldots$。说它们“趋于”$f$，可以只要求每个固定输入处的输出趋于$f$的值，也可以要求同一个序号对全部输入都达到给定精度。两种要求的差别决定了极限能否继承连续性。本节设$X$是集合（讨论连续性时是拓扑空间），$(Y,d)$是度量空间，$f_n,f:X\to Y$。先定义两种收敛并说明一致距离的适用范围，再证明一致极限保持连续，最后说明目标空间完备时，一致柯西列必有一致极限。

\subsection{逐点收敛与一致收敛}
\label{sec:mmap-pointwise-uniform}

两种收敛的区别在于“达到精度所需的序号”是否允许随输入改变。

\begin{definition}[逐点收敛与一致收敛]
\label{def:mmap-pointwise-uniform}
若对每个$x\in X$，当$n\to\infty$时都有$d(f_n(x),f(x))\to0$，即对每个$x$和$\varepsilon>0$存在$N$（可依赖于$x$）使$n\geq N$时$d(f_n(x),f(x))<\varepsilon$，则称$(f_n)$\term{逐点收敛}（pointwise convergence）于$f$。若对每个$\varepsilon>0$存在不依赖于$x$的$N$，使$n\geq N$时对所有$x\in X$都有$d(f_n(x),f(x))<\varepsilon$，则称$(f_n)$\term{一致收敛}（uniform convergence）于$f$。
\end{definition}

在$[0,1]$上取$f_n(x)=x^n$。对$x\in[0,1)$有$x^n\to0$，而$f_n(1)=1$，所以逐点极限是$f(x)=0$（$x<1$）、$f(1)=1$。所需序号随$x$变化：要使$x^n<0.01$，$x=0.9$时需要$n\geq44$，因为$\ln0.01/\ln0.9\approx43.7$；$x=0.99$时需要$n\geq459$，因为$\ln0.01/\ln0.99\approx458.2$。越接近$1$，所需序号越大，没有一个$N$对全部$x\in[0,1)$有效：对任何$n$，$\sup_{0\leq x<1}x^n=1$。所以这个映射列不一致收敛。限制到$[0,1/2]$上，$\sup_{0\leq x\leq1/2}x^n=2^{-n}\to0$，就一致收敛于$0$。

一致收敛可以用一个量来概括。令
\begin{equation}
d_\infty(f,g)=\sup_{x\in X}d(f(x),g(x))\in[0,\infty],
\label{eq:mmap-uniform-distance}
\end{equation}
则$f_n$一致收敛于$f$当且仅当$d_\infty(f_n,f)\to0$。但$d_\infty$可能取值$\infty$：在$\mathbb R$上，$f(x)=x$与$g(x)=0$的上确界距离是$\infty$。因此它不能在全体映射上当作通常的度量。

把映射限制在像集有界的映射上，问题就消失了。若$f(X)$含于$Y$的某个球，称$f$为\term{有界映射}（bounded map），记全体有界映射为$B(X,Y)$。对$f,g\in B(X,Y)$，$f(X)\subseteq B(y_1,r_1)$、$g(X)\subseteq B(y_2,r_2)$，则$d(f(x),g(x))<r_1+d(y_1,y_2)+r_2$，所以$d_\infty(f,g)<\infty$。三角不等式逐点成立再取上确界，$d_\infty(f,h)\leq d_\infty(f,g)+d_\infty(g,h)$，零距离判定也逐点成立，所以$d_\infty$是$B(X,Y)$上的度量，称为\term{一致度量}（uniform metric）。在一般映射上谈一致收敛时，仍然使用定义~\ref{def:mmap-pointwise-uniform}的$\varepsilon$-$N$表述，只要求$d_\infty(f_n,f)$最终有限并趋于零。

\subsection{一致极限的连续性}
\label{sec:mmap-uniform-limit-continuity}

$x^n$的例子中，每个$f_n$都连续，逐点极限却在$x=1$处跳跃。一致收敛排除了这种情形，证明只需把误差分成三段。

\begin{theorem}[一致极限保持连续]
\label{thm:mmap-uniform-limit-continuous}
设$X$是拓扑空间，$(Y,d)$是度量空间，$f_n:X\to Y$连续，并且一致收敛于$f$。则$f$连续。
\end{theorem}

\begin{proof}
设$V\subseteq Y$是开集，$x\in f^{-1}(V)$。取$\varepsilon>0$使$B(f(x),\varepsilon)\subseteq V$，再取$N$使对所有$x'\in X$都有$d(f_N(x'),f(x'))<\varepsilon/3$。$f_N$连续，所以$U=f_N^{-1}\bigl(B(f_N(x),\varepsilon/3)\bigr)$是含$x$的开集。对$x'\in U$，
\begin{equation*}
\begin{aligned}
d(f(x'),f(x))
&\leq d(f(x'),f_N(x'))+d(f_N(x'),f_N(x))+d(f_N(x),f(x))\\
&<\tfrac\varepsilon3+\tfrac\varepsilon3+\tfrac\varepsilon3=\varepsilon,
\end{aligned}
\end{equation*}
所以$U\subseteq f^{-1}(V)$。$f^{-1}(V)$的每个点都有含于其中的开集，它是开集。
\end{proof}

三段误差中，首尾两段由一致收敛控制，并且用的是同一个$N$；中间一段由单个$f_N$的连续性控制。若只有逐点收敛，首段$d(f(x'),f_N(x'))$所需的$N$随$x'$变化，无法与中间一段的邻域$U$同时取定，$x^n$正是这样失败的。

定理的条件很少：定义域只需是拓扑空间，不要求紧致；目标空间只需是度量空间，不要求完备，因为极限$f$已经假定存在。一致收敛也只是充分条件：$x^n$在$[0,1)$上逐点收敛于连续函数$0$，但不一致收敛。

\subsection{一致柯西列的极限}
\label{sec:mmap-uniform-cauchy}

上一小节假定了极限$f$存在。和点列一样，实际中往往只能看到映射列的后项彼此接近，需要从这一点推出极限存在。

\begin{definition}[一致柯西列]
\label{def:mmap-uniform-cauchy}
若对任意$\varepsilon>0$，存在$N$，使$m,n\geq N$时对所有$x\in X$都有$d(f_m(x),f_n(x))<\varepsilon$，则称$(f_n)$为\term{一致柯西列}（uniformly Cauchy sequence）。
\end{definition}

\begin{theorem}[一致柯西列的一致极限]
\label{thm:mmap-uniform-cauchy-limit}
设$(Y,d)$完备，$(f_n)$是$X\to Y$的一致柯西列。则存在$f:X\to Y$，使$f_n$一致收敛于$f$。
\end{theorem}

\begin{proof}
对每个固定的$x$，$(f_n(x))$是$Y$中的柯西列，由完备性收敛，记极限为$f(x)$。给定$\varepsilon>0$，取$N$使$m,n\geq N$时对所有$x$有$d(f_n(x),f_m(x))<\varepsilon$。固定$n\geq N$和$x$，令$m\to\infty$。由第~\ref{chap:metric-spaces}章的不等式$\lvert d(a,b)-d(a,c)\rvert\leq d(b,c)$，$d(f_n(x),f_m(x))\to d(f_n(x),f(x))$，所以$d(f_n(x),f(x))\leq\varepsilon$。这对所有$x$和所有$n\geq N$成立，即一致收敛。
\end{proof}

证明的顺序值得注意：先逐点取极限得到候选的$f$，再在不等式中让后项趋于极限，把对所有$x$统一的估计传给$f$。目标空间的完备性不可省略。取$Y=\mathbb Q$，$f_n$为常值映射，取值为$\sqrt2$的第$n$个十进制截断，它们是一致柯西列，却没有取值在$\mathbb Q$中的极限。

把本小节与上一小节合起来，就得到函数空间完备性的共同依据。

\begin{theorem}[有界映射空间的完备性]
\label{thm:mmap-bounded-maps-complete}
设$(Y,d)$完备。则$(B(X,Y),d_\infty)$完备。若$X$是拓扑空间，有界连续映射全体$C_b(X,Y)$是$B(X,Y)$的闭子集，因而也完备。
\end{theorem}

\begin{proof}
$B(X,Y)$中关于$d_\infty$的柯西列就是一致柯西列，由定理~\ref{thm:mmap-uniform-cauchy-limit}一致收敛于某个$f$。取$N$使$d_\infty(f_N,f)\leq1$，$f_N$的像含于某个球$B(y,r)$，则$f$的像含于$B(y,r+1)$，$f\in B(X,Y)$，并且$d_\infty(f_n,f)\to0$。若各$f_n$还连续，由定理~\ref{thm:mmap-uniform-limit-continuous}，$f$连续。这说明$C_b(X,Y)$包含其中点列的极限，是闭集；按第~\ref{chap:metric-spaces}章，完备空间的闭子空间完备。
\end{proof}

$X$紧致时，连续映射的像是紧集，因而有界，所以$C_b(X,Y)$就是全体连续映射$C(X,Y)$。例如$C([0,1],\mathbb R)$在一致度量下完备。检验具体的函数列时，常常先估计一致柯西性。在$[0,1]$上令$s_n(x)=\sum_{k=1}^nx^k/k^2$。对$m>n$，
\begin{equation*}
\lvert s_m(x)-s_n(x)\rvert\leq\sum_{k=n+1}^m\frac1{k^2}\leq\sum_{k=n+1}^m\Bigl(\frac1{k-1}-\frac1k\Bigr)<\frac1n,
\end{equation*}
上界与$x$无关，所以$(s_n)$是一致柯西列，一致收敛于一个连续函数。第~\ref{chap:normed-spaces}章核验具体函数空间的完备性时，反复使用的就是这一模式：先逐点取极限，再证明极限属于所考察的空间。
